% ============ 05. Calidad de datos y gobernanza ============
\section{Calidad de datos, gobernanza y reproducibilidad para observatorios}

\subsection{Qué dice la literatura}

La gobernanza de datos para observatorios de investigación se apoya en tres
pilares consolidados: (i) las \textbf{dimensiones de calidad de datos} de
\citet{wang1996beyond} (exactitud, completitud, consistencia, oportunidad,
unicidad, validez) y su operacionalización en marcos de gestión (DAMA-DMBOK,
\citet{dama2017dmbok}; ISO/IEC 25012); (ii) los \textbf{principios FAIR}
\citep{wilkinson2016fair} (Findable, Accessible, Interoperable, Reusable) como
estándar para datos de investigación; y (iii) la \textbf{reproducibilidad
técnica}: contenedores, gestión de versiones de datos (\texttt{DVC}, Git LFS),
gestores de flujos de trabajo (Nextflow, Snakemake) y \emph{data observability}
— monitoreo de la salud de los datos dentro del pipeline
\citep{souza2023observability} (preprint IEEE eScience 2023, sección 13.2). La
literatura reciente enfatiza el \textbf{linaje/provenance} (de dónde vino cada
valor y qué transformaciones sufrió) como condición para auditar resultados.

\subsection{Relación con este repositorio}

La auditoría del proyecto (documento maestro) verificó: catálogo YAML con
diccionario de datos y decisiones documentadas (buena base FAIR); pero
dimensiones de calidad evaluadas como parciales (exactitud sin contraste
externo; consistencia afectada por mojibake y duplicidad de artefactos;
oportunidad con pull 2021--2026); datos crudos y procesados sin versionar
(ausencia de DVC pese a la nota del \texttt{.gitignore}); CI/CD rota
(módulos inexistentes en el workflow); contrato de esquema (\texttt{pandera})
declarado y sin uso; y ausencia de linaje explícito entre ingesta,
transformación y evidencias.

\subsection{Mejoras recomendadas}

\begin{enumerate}
  \item Formalizar una \textbf{auditoría de calidad por las 6 dimensiones} de
        \citet{wang1996beyond} con reportes automáticos por convocatoria, y
        declarar el estado por dimensión (completitud, unicidad, validez,
        consistencia, exactitud, oportunidad).
  \item Implementar \textbf{FAIR práctico}: documentando en el catálogo con
        identificadores persistentes, contratos de esquema
        (\texttt{pandera}/\texttt{great\_expectations}) en ingesta y
        transformación, y un registro de \textbf{linaje} (archivo→transformación
        →evidencia).
  \item Versionar datos críticos con \textbf{DVC o Git LFS} (incluido el
        consolidado completo de 6 convocatorias, hoy ausente del clon) y
        reparar CI/CD para que la ingesta periódica sea realmente reproducible
        (recomendación de \citep{souza2023observability} sobre observabilidad
        en pipelines). La arquitectura dimensional del repositorio (esquema
        estrella en DuckDB) sigue el modelo de \citet{kimball2013data},
        apropiado para consultas analíticas; la literatura recomienda
        complementarlo con linaje y validación de contratos para convertirlo en
        un observatorio auditable.
\end{enumerate}
